\section{What experts identify at the institutional level}
Experts distinguished formal recognition of a document from successful verification at a provider. They raised continuity when documentation changes, and practical disruptions such as limited data, unstable phone numbers and difficulties with app authentication. They proposed human assistance at onboarding and first use alongside simpler interfaces. This is an implementation diagnosis, not an estimate of the effect of training or support.

They also distinguished domestic payment interoperability from the cross-border work of verifying identity, converting currency, delivering funds at destination and handling disputes across organizations. A faster domestic payment can improve one leg without changing the others. In hybrid transactions the protection attached to that leg may not settle who is responsible for the rest of the transfer. Experts identified separate responsibilities for migration authorities, regulators, payment operators, providers and organizations offering assistance. The issue they raised was coordination across those responsibilities, not simply whether a rail or complaint channel existed. For a domestic Bre-B transaction, for example, the central bank directs customers to the participating financial institution's complaint channel \citep{banrepfaq}. That documented route does not say what occurred in any interviewee's case or establish a remedy for a cross-border dispute.
